\subsubsection{Diagrama de árbol de problemas}

En los centros de salud de San Marcos, la gestión de turnos presenta problemas como largas filas, desorganización en la atención y falta de prioridad para pacientes con emergencias, adultos mayores y mujeres embarazadas. Esto genera retrasos, confusión e insatisfacción tanto en los pacientes como en el personal de salud. Para solucionar esta problemática, se propone implementar un sistema automatizado de gestión de turnos que permita registrar pacientes digitalmente, asignar turnos según prioridad y mejorar la organización y eficiencia de la atención médica.

Problemas en la Gestión de Turnos en Centros de Salud de San Marcos, Carazo

PROBLEMA CENTRAL

Problemas en la Gestión de Turnos

\paragraph{Causas}

Turnos manuales

Falta de organización

Demora en la atención

\paragraph{Consecuencias}

Largas esperas

Pacientes molestos

Errores en turnos
